\section{A proof-bearing calculus of representation change}
\label{sec:calculus}

The calculus is intentionally small.  It does not attempt to encode every
feature of a physical theory.  It records the objects needed to decide whether
a particular claim can cross a particular interface.

\subsection{Representation packages}

\begin{definition}[Representation package]
\label{def:package}
A \emph{representation package} is a tuple
\[
 \Rep=(\mathcal K,\mathcal G,\Aphys,\mathcal S,\mathcal T,
       \mathcal L,\mathcal Q)
\]
with the following typed components.
\begin{enumerate}[label=(\roman*),itemsep=2pt]
 \item $\mathcal G\rightrightarrows\mathcal K$ is a groupoid of
 kinematical presentations and redundancies.
 \item $\Aphys$ is a unital $*$-algebra of physical observables, equipped
 with seminorms $p_\alpha$ and an interpretation of support.
 \item $\mathcal S\subset\Aphys^{*}$ is a specified class of normalized
 positive states.
 \item $\mathcal T$ is physical time evolution, a Euclidean correlation
 family, or another explicitly typed dynamics; these alternatives are not
 identified.
 \item $\mathcal L=(I,\preceq,\lambda,\mathfrak C)$ records the directed
 index set, physical scales and units, admissible limiting paths, and
 convergence topologies.
 \item $\mathcal Q$ is a family of property predicates, each with a bearer,
 domain, units, and quantifier prefix.
\end{enumerate}
The algebra $\Aphys$ is invariant under $\mathcal G$, or is supplied with a
descent theorem from a covariant algebra.
\end{definition}

The groupoid accommodates the fact that a gauge transformation relates
presentations rather than physical histories.  The distinction between
$\mathcal T$ and $\mathcal L$ prevents physical time from being identified
with an RG parameter.  The property set $\mathcal Q$ is not a list of words:
``locality,'' for example, must specify a net or support rule, while ``gap''
must specify a generator, a vacuum sector, and physical energy units.
Appendix~\ref{app:typed} gives the full typing and identity conditions.

\begin{example}[Finite semigroup package]
Let $\mathcal K=\mathbb C^{d}$, let the redundancy groupoid be trivial, let
$\Aphys=M_d(\mathbb C)$, and let $\mathcal S$ be density matrices.  A positive
self-adjoint $H$ defines $\mathcal T_t(A)=\e^{-tH}A\e^{-tH}$ for Euclidean
quadratic forms.  The limit data may be a sequence of inclusions into a
separable Hilbert space.  Spectral support, core completeness, and semigroup
decay are well-typed property predicates.
\end{example}

\begin{example}[Regulated lattice gauge package]
For a finite lattice, $\mathcal K=G^{E}$, the vertex gauge group acts as a
groupoid of link presentations, and $\Aphys$ is a chosen algebra of invariant
functions.  A positive lattice measure defines states and Euclidean
correlations.  The limit data include lattice spacing, physical volume,
bare/renormalized parameters, and observable renormalization.  Gauge
invariance is exact at finite cutoff, while continuum covariance, physical
observable completeness, and a Hamiltonian gap are limit or reconstruction
predicates.  This instantiates every component without treating the regulator
as the continuum theory.
\end{example}

\subsection{Certificate-profiled morphisms}

Let $F:\Rep_0\to\Rep_1$ contain a pullback
$F^*: \mathcal D_1\to\Aphys^{(0)}$ on a declared target-observable domain and,
when defined, a state pushforward $F_*:\mathcal S_0\to\mathcal S_1$.  Exact
duality would require
\[
 \langle F_*\omega,A\rangle
   =\langle\omega,F^*A\rangle.
\]
Approximate, limiting, and coarse-graining maps instead carry a family
$\varepsilon_F(A;i)$ satisfying
\begin{equation}
\label{eq:fidelity}
 \left|\langle F_*\omega,A\rangle-
        \langle\omega,F^*A\rangle\right|
 \leq \varepsilon_F(A;i)
\end{equation}
for the stated states, observables, and indices.

\begin{definition}[Certificate profile]
\label{def:profile}
The profile of $F$ is
\[
 \Cert(F)=(\Exact_F,\varepsilon_F,\Forgot_F,\Recover_F,\Uniform_F).
\]
$\Exact_F$ contains exact property channels.  $\varepsilon_F$ is a
heterogeneous family of observable-dependent moduli.  $\Forgot_F$ lists
structures not available in the target variables.  $\Recover_F$ contains
theorems or partial sections that recover specified claims despite such loss.
$\Uniform_F$ records the full index domain and order of quantifiers on which
each entry holds.
\end{definition}

Forgetting is therefore not the negation of physical existence.  It says that
the target representation does not by itself carry a datum required by the
inference.  Recovery is not a declaration that the datum ``must still be
there.''  It names a map or theorem, its domain, and a modulus.  Exact
inverses, reconstruction on an observable core, asymptotic recovery, and
external reconstruction theorems are consequently different entries.

Figure~\ref{fig:profile} displays the logical form.  The arrow can preserve
some channels and erase others; a recovery theorem is a second arrow, not an
optimistic interpretation of the first.

\begin{figure}[t]
\centering
\begin{tikzpicture}[
  box/.style={draw,rounded corners,align=center,minimum width=34mm,minimum height=10mm},
  cert/.style={draw,dashed,rounded corners,align=left,inner sep=5pt},
  arr/.style={-{Latex[length=2.2mm]},thick},
  node distance=18mm and 25mm]
\node[box] (r0) {$\Rep_0$\\source description};
\node[box,right=of r0] (r1) {$\Rep_1$\\target description};
\draw[arr] (r0) -- node[above] {$F$} (r1);
\node[cert,below=of r0] (c)
 {$\Exact_F$: exact channels\\
  $\varepsilon_F$: quantitative fidelity\\
  $\Forgot_F$: unavailable structure\\
  $\Uniform_F$: index domain};
\node[box,below=of r1] (p) {target property\\$Q_1$};
\draw[arr,dashed] (r1) -- node[right] {direct} (p);
\draw[arr,blue!65!black] (c.east) --
 node[below,align=center,font=\small] {$\Recover_F$} (p.west);
\end{tikzpicture}
\caption{A property-relative representation morphism.  Recovery is a
separate theorem-backed route from the source data and certificate profile to
the target claim.}
\label{fig:profile}
\end{figure}

\subsection{Composition without scalarization}

Suppose $F:\Rep_0\to\Rep_1$ and $G:\Rep_1\to\Rep_2$.  For a target
observable $B$, inserting the intermediate expectation gives
\begin{equation}
\label{eq:composition}
 \varepsilon_{G\circ F}(B;i)
 \leq \varepsilon_G(B;i)+\varepsilon_F(G^*B;i).
\end{equation}
If the pullback of $G$ obeys
$p_1(G^*B)\leq L_G(B)p_2(B)$ and the first error is controlled by $p_1$,
then
\[
 \varepsilon_{G\circ F}(B;i)
 \leq \varepsilon_G(B;i)+L_G(B)\varepsilon_F(B;i).
\]
Locality range, positivity, gauge descent, and spectral fidelity are not
folded into this scalar expression.  Each has its own composition rule and
failure mode.

\begin{proposition}[Typed profile composition]
\label{prop:composition}
On composable observable domains, expectation moduli compose by
Equation~\eqref{eq:composition}; exact channels compose by conjunction;
uniformity domains compose by intersection with the pulled-back quantifier
order; and a forgotten prerequisite can appear in the composite conclusion
only through an explicit recovery theorem.  For three composable morphisms,
the expectation modulus is associative:
\[
 \varepsilon_H+L_H\varepsilon_G+L_HL_G\varepsilon_F,
\]
independently of bracketing.
\end{proposition}

The proof is in Appendix~\ref{app:composition}.  This proposition is still
bookkeeping.  Its importance is negative and architectural: it prevents one
successful channel from licensing another.  The substantive result comes
when a recovery theorem turns a family of Euclidean bounds into a spectral
property.

\subsection{How much category theory is claimed?}

Exact representation equivalences, with total observable maps and identities,
form an ordinary category.  Approximate arrows with seminorm-valued profiles
can be organized in enriched or ordered structures, and institution-style
semantics provides useful analogies for invariance under change of notation
\cite{GoguenBurstall1992}.  We do not claim that every coarse graining,
limiting correspondence, or partial reconstruction belongs to one completed
category.  In the present paper, ``morphism'' means a typed composable arrow
on the stated domain.  Categorical vocabulary is retained only where identity
and associativity have been proved.
